\documentclass[UTF8,zihao=-4]{ctexart}
\usepackage[a4paper,margin=1.85cm]{geometry}
\usepackage{hyperref}
\usepackage{graphicx}
\usepackage{verbatim}
\usepackage{xcolor}
\usepackage{fancyhdr}
\hypersetup{hidelinks}
\setlength{\parindent}{0pt}
\setlength{\parskip}{0.45em}
\setlength{\headheight}{14.5pt}
\emergencystretch=4em
\sloppy
\pagestyle{fancy}
\fancyhf{}
\lhead{编译原理完整模拟卷}
\rhead{\thepage}
\title{Slides 程序分析完整模拟卷 09（题目与完整解析）}
\author{}
\date{}
\begin{document}
\maketitle
\setcounter{tocdepth}{1}
\tableofcontents
\newpage

\section*{卷面说明}
\addcontentsline{toc}{section}{卷面说明}


\begin{quote}
\footnotesize
\begin{verbatim}
考试时间：120 分钟
总分：100 分
说明：本卷独立覆盖流敏感、路径敏感、上下文敏感、符号执行、SAT、CNF、Tseitin 和可满足性等价证明。
\end{verbatim}
\end{quote}

\section*{A 卷题目}
\addcontentsline{toc}{section}{A 卷题目}


\subsection*{一、三种敏感性与符号状态（20 分）}
\addcontentsline{toc}{subsection}{一、三种敏感性与符号状态（20 分）}


给定程序：

\begin{quote}
\footnotesize
\begin{verbatim}
x = 0;
if (c)
    x = 1;
y = x;

int id(int p) { return p; }
a = id(1);      // call site C1
b = id(2);      // call site C2
\end{verbatim}
\end{quote}

\noindent 1. 分别定义 flow-sensitive、path-sensitive、context-sensitive。\\
\noindent 2. 流不敏感分析把 \texttt{x=0} 和 \texttt{x=1} 合并后，在 \texttt{y=x} 处得到什么集合。\\
\noindent 3. 流敏感但路径不敏感分析在 \texttt{y=x} 处得到什么集合。\\
\noindent 4. 路径敏感分析在 \texttt{y=x} 处保留哪两组“路径条件 + x 值”。\\
\noindent 5. 上下文不敏感和上下文敏感分析 \texttt{id} 时，\texttt{a、b} 的返回值结果分别是什么。\\

\subsection*{二、符号执行路径枚举（20 分）}
\addcontentsline{toc}{subsection}{二、符号执行路径枚举（20 分）}


程序：

\begin{quote}
\footnotesize
\begin{verbatim}
if (x > 0)
  y = x + 1;
else
  y = 1 - x;

if (y == 2)
  error();
\end{verbatim}
\end{quote}

\noindent 1. 写出第一条路径的路径条件和符号状态。\\
\noindent 2. 写出第二条路径的路径条件和符号状态。\\
\noindent 3. 判断 \texttt{error()} 是否可达。\\
\noindent 4. 给出能到达 \texttt{error()} 的输入。\\
\noindent 5. 说明符号执行为什么可能发生路径爆炸。\\

\subsection*{三、SAT、Validity 与 CNF（20 分）}
\addcontentsline{toc}{subsection}{三、SAT、Validity 与 CNF（20 分）}


回答下列问题：

\noindent 1. 什么是 satisfiable。\\
\noindent 2. 什么是 valid。\\
\noindent 3. satisfiability 和 validity 的对偶关系是什么。\\
\noindent 4. 什么是 CNF。\\
\noindent 5. 判断下列表达式是否为 CNF，并说明原因：\\

\begin{quote}
\footnotesize
\begin{verbatim}
(a or b) and (not c or d)
not (a and b)
(a or (b and c))
\end{verbatim}
\end{quote}

\subsection*{四、Tseitin Transformation（20 分）}
\addcontentsline{toc}{subsection}{四、Tseitin Transformation（20 分）}


对公式：

\begin{quote}
\footnotesize
\begin{verbatim}
F = (a or b) and not (c and d)
\end{verbatim}
\end{quote}

\noindent 1. 为子公式引入新变量。\\
\noindent 2. 写出每个新变量与子公式等价的约束。\\
\noindent 3. 把约束转换成 CNF 子句。\\
\noindent 4. 写出最终的 equisatisfiable CNF。\\
\noindent 5. 说明 Tseitin 转换为什么比直接展开更适合 SAT 求解。\\

\subsection*{五、Equisatisfiability 证明（20 分）}
\addcontentsline{toc}{subsection}{五、Equisatisfiability 证明（20 分）}


证明下列两个公式 equisatisfiable：

\begin{quote}
\footnotesize
\begin{verbatim}
F1 = (a0 or a1 or ... or an or c) and (b0 or b1 or ... or bm or not c)
F2 = (a0 or a1 or ... or an or b0 or b1 or ... or bm)
\end{verbatim}
\end{quote}

\noindent 1. 设 \texttt{A = a0 or a1 or ... or an}，\texttt{B = b0 or b1 or ... or bm}，重写两个公式。\\
\noindent 2. 证明若 \texttt{F1} 可满足，则 \texttt{F2} 可满足。\\
\noindent 3. 证明若 \texttt{F2} 可满足，则存在 \texttt{c} 的取值使 \texttt{F1} 可满足。\\
\noindent 4. 说明 equisatisfiable 和 equivalent 的区别。\\
\noindent 5. 说明该证明在 Tseitin 转换中的作用。\\

\newpage
\section*{B 完整解析}
\addcontentsline{toc}{section}{B 完整解析}


\subsection*{一、三种敏感性与符号状态解析}
\addcontentsline{toc}{subsection}{一、三种敏感性与符号状态解析}


flow-sensitive：

\begin{quote}
\footnotesize
\begin{verbatim}
分析结果区分程序点，考虑语句执行顺序。
同一变量在不同程序点可以有不同抽象信息。
\end{verbatim}
\end{quote}

path-sensitive：

\begin{quote}
\footnotesize
\begin{verbatim}
分析结果区分不同控制流路径。
同一程序点如果由不同路径到达，可以保留不同路径条件和不同状态。
符号执行是典型的路径敏感分析。
\end{verbatim}
\end{quote}

context-sensitive：

\begin{quote}
\footnotesize
\begin{verbatim}
过程间分析时区分不同调用上下文。
同一函数从不同调用点、不同调用链进入时，可以得到不同分析结果。
\end{verbatim}
\end{quote}

程序分析结果：

\begin{quote}
\footnotesize
\begin{verbatim}
流不敏感：
忽略两条赋值的先后关系，把 x 的可能值合并为 {0,1}，所以 y in {0,1}。

流敏感但路径不敏感：
分析知道 x=0 先执行，再经过条件分支；在合流点把两条路径合并，仍得到 y in {0,1}。

路径敏感：
路径 c=true：x=1，y=1。
路径 c=false：x=0，y=0。
两组状态分别保留，不立即合并。
\end{verbatim}
\end{quote}

调用上下文结果：

\begin{quote}
\footnotesize
\begin{verbatim}
上下文不敏感：
把 C1、C2 两次调用合并分析，id 的参数/返回值集合为 {1,2}，因此 a、b 都被近似为 {1,2}。

上下文敏感：
区分 C1、C2，得到 a=1、b=2。
\end{verbatim}
\end{quote}

\subsection*{二、符号执行路径枚举答案}
\addcontentsline{toc}{subsection}{二、符号执行路径枚举答案}


设输入符号为：

\begin{quote}
\footnotesize
\begin{verbatim}
sigma0(x) = alpha
\end{verbatim}
\end{quote}

路径一：进入 then 分支：

\begin{quote}
\footnotesize
\begin{verbatim}
条件：alpha > 0
执行：y = x + 1
状态：sigma(y) = alpha + 1

若进入 error：
路径条件 pi = alpha > 0 and alpha + 1 == 2
化简得到 alpha = 1。
该路径可满足。
\end{verbatim}
\end{quote}

路径二：进入 else 分支：

路径枚举表：
\begin{center}\footnotesize
\begin{tabular}{llll}
路径 & 路径条件 & error 条件 & 可满足输入\\
then & alpha > 0 & alpha + 1 == 2 & alpha = 1\\
else & alpha <= 0 & 1 - alpha == 2 & alpha = -1
\end{tabular}
\end{center}


\begin{quote}
\footnotesize
\begin{verbatim}
条件：alpha <= 0
执行：y = 1 - x
状态：sigma(y) = 1 - alpha

若进入 error：
路径条件 pi = alpha <= 0 and 1 - alpha == 2
化简得到 alpha = -1。
该路径可满足。
\end{verbatim}
\end{quote}

\texttt{error()} 可达：

\begin{quote}
\footnotesize
\begin{verbatim}
可达。
输入 x = 1 走 then 分支并使 y = 2。
输入 x = -1 走 else 分支并使 y = 2。
\end{verbatim}
\end{quote}

路径爆炸：

\begin{quote}
\footnotesize
\begin{verbatim}
每个条件分支都可能把路径数乘 2。
循环和递归还会产生大量甚至无限路径。
因此符号执行通常需要路径裁剪、循环界限、状态合并或约束求解优化。
\end{verbatim}
\end{quote}

\subsection*{三、SAT、Validity 与 CNF 答案}
\addcontentsline{toc}{subsection}{三、SAT、Validity 与 CNF 答案}


satisfiable：

\begin{quote}
\footnotesize
\begin{verbatim}
公式可满足，表示存在一个解释或赋值，使公式取真。
\end{verbatim}
\end{quote}

valid：

\begin{quote}
\footnotesize
\begin{verbatim}
公式有效，表示所有解释或赋值下公式都取真。
\end{verbatim}
\end{quote}

对偶关系：

\begin{quote}
\footnotesize
\begin{verbatim}
公式 F valid，当且仅当 not F unsatisfiable。
公式 F satisfiable，当且仅当 not F not valid。
\end{verbatim}
\end{quote}

CNF：

\begin{quote}
\footnotesize
\begin{verbatim}
CNF 是合取范式。
它是若干子句的合取，每个子句是若干文字的析取。
文字是变量或变量的否定。
\end{verbatim}
\end{quote}

判断：

\begin{quote}
\footnotesize
\begin{verbatim}
(a or b) and (not c or d)
是 CNF。它是两个析取子句的合取。

not (a and b)
不是直接的 CNF，因为否定作用在复合公式上。
可等价改写为 not a or not b，改写后是一个 CNF 子句。

(a or (b and c))
不是 CNF，因为析取内部嵌套了合取。
可分配为 (a or b) and (a or c)。
\end{verbatim}
\end{quote}

\subsection*{四、Tseitin Transformation 答案}
\addcontentsline{toc}{subsection}{四、Tseitin Transformation 答案}


引入变量：

Tseitin 变量表：
\begin{center}\footnotesize
\begin{tabular}{ll}
新变量 & 子公式\\
x1 & a or b\\
x2 & c and d\\
x3 & not x2\\
x4 & x1 and x3
\end{tabular}
\end{center}


\begin{quote}
\footnotesize
\begin{verbatim}
x1 <-> (a or b)
x2 <-> (c and d)
x3 <-> not x2
x4 <-> (x1 and x3)
最终要求 x4 为真。
\end{verbatim}
\end{quote}

等价约束转 CNF：

\begin{quote}
\footnotesize
\begin{verbatim}
x1 <-> (a or b):
(not x1 or a or b)
(x1 or not a)
(x1 or not b)

x2 <-> (c and d):
(not x2 or c)
(not x2 or d)
(x2 or not c or not d)

x3 <-> not x2:
(not x3 or not x2)
(x3 or x2)

x4 <-> (x1 and x3):
(not x4 or x1)
(not x4 or x3)
(x4 or not x1 or not x3)

要求 x4：
(x4)
\end{verbatim}
\end{quote}

最终 equisatisfiable CNF：

\begin{quote}
\footnotesize
\begin{verbatim}
(not x1 or a or b)
and (x1 or not a)
and (x1 or not b)
and (not x2 or c)
and (not x2 or d)
and (x2 or not c or not d)
and (not x3 or not x2)
and (x3 or x2)
and (not x4 or x1)
and (not x4 or x3)
and (x4 or not x1 or not x3)
and (x4)
\end{verbatim}
\end{quote}

优势：

\begin{quote}
\footnotesize
\begin{verbatim}
直接等价展开到 CNF 可能指数膨胀。
Tseitin 转换为子公式引入新变量，只要求可满足性等价，不要求完全等价。
因此 CNF 大小通常线性增长，更适合 SAT 求解器。
\end{verbatim}
\end{quote}

\subsection*{五、Equisatisfiability 证明答案}
\addcontentsline{toc}{subsection}{五、Equisatisfiability 证明答案}


重写：

\begin{quote}
\footnotesize
\begin{verbatim}
A = a0 or a1 or ... or an
B = b0 or b1 or ... or bm

F1 = (A or c) and (B or not c)
F2 = A or B
\end{verbatim}
\end{quote}

\texttt{F1 -> F2}：

\begin{quote}
\footnotesize
\begin{verbatim}
若 F1 可满足，则在某个赋值下 (A or c) 为真且 (B or not c) 为真。

如果 c 为真，则第二个子句要求 B 为真，因此 A or B 为真。
如果 c 为假，则第一个子句要求 A 为真，因此 A or B 为真。

所以 F2 可满足。
\end{verbatim}
\end{quote}

\texttt{F2 -> F1}：

\begin{quote}
\footnotesize
\begin{verbatim}
若 F2 可满足，则在某个赋值下 A or B 为真。

如果 A 为真，令 c = false。
则 A or c 为真，B or not c 也因为 not c 为真而为真。

如果 A 为假，则 B 必为真，令 c = true。
则 A or c 为真，B or not c 也因为 B 为真而为真。

因此存在 c 的取值使 F1 可满足。
\end{verbatim}
\end{quote}

equisatisfiable 与 equivalent：

\begin{quote}
\footnotesize
\begin{verbatim}
equisatisfiable 只要求两个公式同时可满足或同时不可满足。
equivalent 要求在所有赋值下真值完全相同。
Tseitin 转换通常只保持可满足性等价，不保持原变量和新变量扩展空间上的完全等价。
\end{verbatim}
\end{quote}

作用：

\begin{quote}
\footnotesize
\begin{verbatim}
Tseitin 转换把复杂公式转为 CNF 后交给 SAT 求解器。
SAT 求解只关心是否存在满足赋值，因此保持 equisatisfiability 就足够。
这也是引入新变量但仍能正确判断原公式可满足性的原因。
\end{verbatim}
\end{quote}
\end{document}
